\printedpage{59}
% head: A. Grothendieck : Dérivateurs, Ch. VII, Édition des Universités de Montpellier II et Paris VII
% foot: Version 22 décembre 2001

\begin{tabular}{lll}
Thus & $1\longmapsto 2$ & corresponds to \quad $F_0$\\[1ex]
 & $1\longmapsto 1$ & corresponds to \quad $F_1$\\[1ex]
 & $1\longmapsto 0$ & corresponds to \quad $F_2$
\end{tabular}

and of course we have $F_2\longrightarrow F_1\longrightarrow F_0$, % [sic] bien sur -> bien sûr, accent missing
\[
\uHom(\Delta^2_\varepsilon,\Delta^2_\varepsilon)\simeq\boldsymbol{\Delta}^2,\qquad\quad\text{0-connected category.}
\]

In these two cases, two elements of $\mathrm{Hom}(\tau,\tau')$ are joined by a homotopy which is \textit{simple.}

\msp{page 93}

B) Case $\tau=\Delta_\varepsilon\Delta^l_{-\varepsilon}\Delta_\varepsilon$, $\tau'=\Delta^2_\varepsilon$ or $\Delta_\varepsilon$.

Let us note that we have, for arbitrary $\tau,\tau'\in\uTy$, an isomorphism of categories
\[
\uHom(\tau,\tau')\longrightarrow\uHom(\tau'^o,\tau^o),
\]
simply expressing that
\[
\tau\longmapsto\tau^o
\]
is a \texttt{[2]}-functor from $\uTy$ to $\uTy$, and even an automorphism of this \texttt{[2]}-category. Thus to say that $\uHom(\tau,\tau')$ is connected amounts to saying that $\uHom(\tau^o,\tau'^o)$ is. Thus we can suppose above that $\varepsilon=+1$ (and the same in A!).
\[
\begin{aligned}
\tau :\quad & \longrightarrow\ \underbrace{\longleftarrow\ \longleftarrow\ \cdots\ \longleftarrow}_{l}\ \longrightarrow\\
\tau' :\quad & \longrightarrow\ \longrightarrow\quad\text{or}\quad\longrightarrow\ .
\end{aligned}
\]

But all the morphisms $\tau\longrightarrow\tau'$ are trivial on the critical interval $I_2$ (retrograde), which implies that
\[
\uHom(\tau,\tau')\xleftarrow{\ \sim\ }\uHom(\Delta^2_\varepsilon,\tau'),
\]
via the canonical contraction functor
\[
I_\tau\longrightarrow I_{\Delta^2_\varepsilon}=\boldsymbol{\Delta}^2.
\]

Now in A we saw that $\uHom(\Delta^2_\varepsilon,\tau')$ is connected (for $\tau'$ of the form $\Delta_\varepsilon$ or $\Delta^2_\varepsilon$).

We've won!

\newpage
\printedpage{60}
% head: A. Grothendieck : Dérivateurs, Ch. VII, Édition des Universités de Montpellier II et Paris VII
% foot: Version 22 décembre 2001

\msp{page 94}

\section*{9\quad Homotopy theorem for categories of paths}

Let $X$ be a category, $\Sigma\subseteq\mathrm{Fl}(X)$ \texttt{[an]} admissible set of arrows in $X$, $a,b\in\mathrm{Ob}\,X$. Let us consider for variable $\tau\in\mathrm{Ty}$, the categories
\[
\uCh_{\Sigma,\tau}(X),\qquad\quad\uCh_{\Sigma,\tau}(X;a,b),
\]
which form the fibers of the Cat-fibered categories
\[
\begin{tikzcd}[column sep=small]
\uCh_\Sigma(X) \arrow[dr] & & \uCh_\Sigma(X;a,b) \arrow[ll, "\text{incl.}"'] \arrow[dl]\\
 & \uTy &
\end{tikzcd}
\]
over the category $\uTy$. Thus, for $\tau,\tau'\in\mathrm{Ty}$, and for any
\[
\varphi:\tau\longrightarrow\tau',
\]
we have a commutative square of homomorphisms
\[
\begin{tikzcd}
\uCh_{\Sigma,\tau'}(X) \arrow[d, "\varphi^*_X"'] & \uCh_{\Sigma,\tau'}(X;a,b) \arrow[l, "\text{incl.}"'] \arrow[d, "\varphi^*_{a,b}"]\\
\uCh_{\Sigma,\tau}(X) & \uCh_{\Sigma,\tau}(X;a,b), \arrow[l, "\text{incl.}"']
\end{tikzcd}
\]
with transitivity for variable $\varphi$ (composition $\psi\circ\varphi$ of arrows in $\uTy$).

In fact, for $\varphi$ variable in $\uHom(\tau,\tau')$, the inverse image functors $\varphi^*_X$, $\varphi^*_{a,b}$ depend \textit{functorially} on $\varphi$; i.e. we have \textit{functors}
\[
\begin{tikzcd}[row sep=small]
 & \uHom(\uCh_{\Sigma,\tau'}(X),\uCh_{\Sigma,\tau}(X))\\
\uHom(\tau,\tau') \arrow[ur] \arrow[dr] & \\
 & \uHom(\uCh_{\Sigma,\tau'}(X;a,b),\uCh_{\Sigma,\tau}(X;a,b))
\end{tikzcd}
\tag*{$(*)$}
\]

\newpage
\printedpage{61}
% head: A. Grothendieck : Dérivateurs, Ch. VII, Édition des Universités de Montpellier II et Paris VII
% foot: Version 22 décembre 2001

\msp{page 95}

This comes from the inclusion
\[
\uHom(\tau,\tau')\hookrightarrow\uHom^{\varepsilon}(I_\tau,I_{\tau'})
\]
and from the description of the categories $\uCh_\tau$, $\uCh_{\tau'}$ as categories of functors, such as
\[
\uCh_{\Sigma,\tau}(X)\simeq\uHom_{\Sigma_\tau,\Sigma}(I_\tau,X),
\]
where $\Sigma_\tau$ denotes the set of arrows in $I_\tau$ which are either invertible or \textit{retrograde.} (This set of arrows is admissible, since stable under composition.) We denote by $\uHom_{\Sigma_\tau,\Sigma}$ the full subcategory of $\uHom$ formed by the functors which map $\Sigma_\tau$ into $\Sigma$. Now in general, if
\[
(I,\Sigma_I),\ \ (I',\Sigma_{I'}),\ \ (X,\Sigma)
\]
are three categories, equipped with admissible sets of arrows, considering the functors $\varphi:I\longrightarrow I'$ with $\varphi(\Sigma_I)\subseteq\Sigma_{I'}$, then the familiar functor
\[
\uHom(I,I')\longrightarrow\uHom(\uHom(I',X),\uHom(I,X))
\]
(transpose of the composition functor
\[
\begin{matrix}
\uHom(I,I') & \times & \uHom(I',X) & \longrightarrow & \uHom(I,X) & \\
(\varphi & , & F') & \longmapsto & F'\circ\varphi & ),
\end{matrix}
\]
generalizes to a functor
\[
\uHom_{\Sigma_I,\Sigma_{I'}}(I,I')\longrightarrow\uHom(\uHom_{\Sigma_{I'},\Sigma}(I',X),\uHom_{\Sigma_I,\Sigma}(I,X)),
\tag*{$(*)$}
\]
transpose of the composition functor
\[
\uHom_{\Sigma_I,\Sigma_{I'}}(I,I')\times\uHom_{\Sigma_{I'},\Sigma}(I',X)\longrightarrow\uHom_{\Sigma_I,\Sigma}(I,X)
\]
induced by $(*)$.

\msp{page 96}

Since every
\[
\varphi:\tau\longrightarrow\tau'
\]
induces
\[
\varphi:I_\tau\longrightarrow I_{\tau'},
\]
transforming retrograde arrows into identity or retrograde arrows, i.e.
\[
\varphi(\Sigma_\tau)\subseteq\Sigma_{\tau'},
\]
\printedpage{62}
% head: A. Grothendieck : Dérivateurs, Ch. VII, Édition des Universités de Montpellier II et Paris VII
% foot: Version 22 décembre 2001
we find the first of the two functors $(*)$ announced on page 94, which is associated by transposition with the functor $\gamma$
\[
\begin{tikzcd}[column sep=large]
\uHom(\tau,\tau')\times\ \underbrace{\uCh_{\Sigma,\tau'}(X)}_{\simeq\,\uHom_{\Sigma_{\tau'},\Sigma}(I_{\tau'},X)} \arrow[r, "\gamma"] \arrow[d, "\text{can.}"'] & \uCh_{\Sigma,\tau}(X) \arrow[d, equal]\\
\uHom_{\Sigma_\tau,\Sigma_{\tau'}}(I_\tau,I_{\tau'})\times\uHom_{\Sigma_{\tau'},\Sigma}(I_{\tau'},X) \arrow[r, dotted, "\gamma_0"] & \uHom_{\Sigma_\tau,\Sigma}(I_\tau,X)
\end{tikzcd}
\tag*{$(*)$}
\]
deduced from the composition functor $\gamma_0$ in this diagram: $\gamma=\gamma_0\circ\mathrm{can}.$ On the other hand, this composition functor $\gamma$ visibly induces a functor on the subcategories of paths with fixed source and target
\[
\uHom(\tau,\tau')\times\uCh_{\Sigma,\tau'}(X;a,b)\longrightarrow\uCh_{\Sigma,\tau}(X;a,b),
\]
whence the second of the two arrows announced in $(*)$ p. 94.

This therefore shows that if
\[
\varphi_1,\varphi_2:\tau\rightrightarrows\tau'
\]
are two morphisms from $\tau$ to $\tau'$ homotopic in $\uHom(\tau,\tau')$, then the induced functors between categories of $\Sigma$-paths or between categories of $\Sigma$-paths with specified endpoints: % [sic] précisés -> précisées, agreement with extrémités

\msp{page 97}
\[
\begin{aligned}
\varphi^*_{1\,X},\varphi^*_{2\,X} : && \uCh_{\Sigma,\tau'}(X) &\quad\rightrightarrows\quad\uCh_{\Sigma,\tau}(X),\\
(\varphi_1)^*_{a,b},(\varphi_2)^*_{a,b} : && \uCh_{\Sigma,\tau'}(X;a,b) &\quad\rightrightarrows\quad\uCh_{\Sigma,\tau}(X;a,b)
\end{aligned}
\]
are also homotopic.

Now the homotopy lemma assures us that the hypothesis $\varphi_1\sim\varphi_2$ is \textit{always} satisfied. Thus:

\textbf{Homotopy theorem on the universal operations on paths.} \textit{Let $X$, $\Sigma$, $a$, $b$ be as above, and $\tau,\tau'\in\mathrm{Ty}$ such that $\mathrm{Hom}(\tau,\tau')\neq\emptyset$. Then the functors $\varphi^*_X$, $\varphi^*_{X;a,b}$ induced by a $\varphi\in\mathrm{Hom}(\tau,\tau')$}
\[
\left\{
\begin{aligned}
\varphi^*_X : && \uCh_{\Sigma,\tau'}(X) &\quad\longrightarrow\quad\uCh_{\Sigma,\tau}(X)\\
\varphi^*_{X;a,b} : && \uCh_{\Sigma,\tau'}(X;a,b) &\quad\longrightarrow\quad\uCh_{\Sigma,\tau}(X;a,b)
\end{aligned}
\right.
\]
\textit{do not depend on $\varphi$,} up to homotopy.

We can make this statement more precise, in such a way that the statement for the $\varphi^*_{X;a,b}$ appears, as it should, as a corollary to the one on $\varphi^*_X$. For this, we are going to look at the categories of paths in $X$ as categories \textit{over $X$,} by the ``source and target'' functors:
\[
\begin{tikzcd}[column sep=small]
\uCh_{\Sigma,\tau}(X) \arrow[dr, "(\sigma_\tau{,}\beta_\tau)"'] & & \uCh_{\Sigma,\tau'}(X) \arrow[dl, "(\sigma_{\tau'}{,}\beta_{\tau'})"]\\
 & X\times X &
\end{tikzcd}
\]

\newpage
\printedpage{63}
% head: A. Grothendieck : Dérivateurs, Ch. VII, Édition des Universités de Montpellier II et Paris VII
% foot: Version 22 décembre 2001

\msp{page 98}

Then the functor $(*)$ (p. 94) factors through the subcategory $\uHom_{\mathrm{Cat}/X\times X}(\cdot,\cdot)$ of $\uHom_{\mathrm{Cat}}(\cdot,\cdot)$ formed by the functors which commute with the projection functors to $X\times X$, and the homomorphisms between such functors which are effected by ``vertical'' arrows of the target category (i.e. by arrows over identity arrows of the base category $X\times X$):
\[
\uHom(\tau,\tau')\longrightarrow\uHom_{\mathrm{Cat}/X\times X}(\uCh_{\Sigma,\tau'}(X),\uCh_{\Sigma,\tau}(X)).
\]

Thus in the statement of the homotopy theorem, we can specify that the homotopy takes place in this subcategory, i.e. is effected by a sequence of simple homotopies which are morphisms of functors between categories over $X\times X$. This then implies the homotopy statement on the $\varphi^*_{X;a,b}$.